\section{Benefícios e Limitações}

A Tevel declara colheita seletiva, preservação da qualidade, ampliação do período de trabalho e redução de custos~\cite{tevel2026}. A disponibilidade operacional e a identificação da maturação favorecem a colheita no momento adequado. Neste estudo, a principal vantagem potencial é ampliar a capacidade de coleta quando falta mão de obra, benefício cuja magnitude exige validação por safra e condição produtiva.

Os registros por fruto e recipiente permitem caracterizar a produção e organizar o processamento posterior. Seu valor depende da confiabilidade, acessibilidade e utilização pelo produtor. O volume de dados não assegura melhores decisões sem calibração das medições e associação correta ao lote.

Como implicação arquitetural, o processamento local reduz a dependência do envio de imagens à infraestrutura remota. A concentração computacional em solo permite empregar processadores mais capazes sem incorporar sua massa aos robôs. O compartilhamento exige avaliar contenção de recursos, prioridades e propagação de falhas aos veículos.

Os cabos condicionam o espaço de trabalho à posição da base e ao percurso entre a vegetação. Devem ser avaliados acesso aos frutos, risco de enrosco, interação entre robôs e mobilidade da plataforma. A documentação consultada não estabelece limites universais de altura, declividade, densidade de copa ou velocidade do vento.

A oclusão limita a observabilidade dos frutos, e a classificação por cor não valida todos os atributos de maturação ou qualidade interna. A aplicação a outra cultivar ou sistema de condução exige avaliar separadamente detecção, acesso, desprendimento e danos. Esses critérios não constituem falhas quantitativamente demonstradas do produto.

A alimentação por cabo transfere a origem da energia de voo, sem eliminar seu consumo. Sem potência e tempo efetivo dos robôs, energia consumida pela plataforma e massa comercializável, não se calcula o consumo específico nem se sustenta superioridade ambiental. A operação prolongada exige inspeção dos cabos, propulsores, dispositivos de coleta e infraestrutura terrestre.

A segurança deve abranger contato com pessoas, falhas de percepção, perda de comunicação local e falhas de alimentação. As proteções de software apresentadas não comprovam, isoladamente, comportamento seguro nessas contingências. A colheita autônoma tampouco dispensa preparação, supervisão, manutenção e logística por operadores.

As fontes consultadas não apresentam preços contratuais, custos completos por quilograma ou retorno do investimento nas implantações analisadas. A viabilidade depende de produtividade útil, utilização anual, transporte, suporte e trabalho complementar. Também exige considerar acesso e exportação dos dados, continuidade das atualizações e dependência do fornecedor. Benefício técnico e viabilidade econômica requerem avaliação conjunta, sem que um garanta o outro.